%% 03_background.tex
%% Section 3: Institutional Background

\section{Institutional Background}
\label{sec:background}

This section describes the three areas of the coal-banking-sovereign nexus in Indonesia. We first outline the structure and significance of Indonesia's coal sector, then characterise the banking sector's exposure to coal, and finally document the sovereign fiscal dependence on coal revenues. The section concludes by reviewing the evolving regulatory landscape that will shape the pace and cost of the transition.

\subsection{Indonesia's Coal Sector}
\label{subsec:coal_sector}

Indonesia is the world's largest exporter of thermal coal, shipping over 500 million tonnes annually and supplying approximately one-quarter of global seaborne trade \citep{IEA2024}. Total domestic production exceeded 600 million tonnes (Mt) in 2023, ranking Indonesia fourth globally after China, India, and the United States in absolute output, but first among net exporters. Coal contributed approximately 10\% of merchandise export revenue and remains the single largest source of primary energy for domestic power generation, fuelling roughly 60\% of the national electricity mix \citep{IESR2023}.

Indonesian coal spans a wide calorific range: 23 of the 34 companies in our sample produce sub-bituminous coal (4,000--5,000 kcal/kg), with the remainder producing bituminous coal (above 6,000 kcal/kg). This distinction matters for stranding analysis because higher-calorie coals command a premium that partially offsets transition-driven price declines.

The sector exhibits moderate concentration: the four largest producers --- AADI (65.8 Mt), BYAN (56.9 Mt), BUMI (43.3 Mt), and the state-owned PTBA (43.3 Mt) --- account for 44\% of listed production. Nearly all mining is open-pit, reflecting shallow coal basins in Sumatra and Kalimantan, yielding a median all-in sustaining cost of approximately \$75/tonne (range: \$10--169/tonne). Mine lives are assumed at 20 years, reflecting the standard IUP concession duration.

Regarding domestic market, a distinctive feature of the Indonesian coal market is the government-mandated Domestic Market Obligation, which requires producers to sell a minimum of 25\% of their planned production to domestic buyers, primarily state electricity utility PT PLN, at a capped price (currently \$70 per tonne for power generation). The DMO compresses margins on domestic sales and effectively creates a cross-subsidy from export earnings to domestic electricity costs. For stranding analysis, the DMO introduces a floor on domestic demand that may slow the pace of coal phase-down but simultaneously reduces per-tonne profitability, making the sector more vulnerable to international price declines \citep{Burke2018, Makhijani2015}.

Indonesian thermal coal is exported primarily to China, India, Japan, South Korea, and Southeast Asian neighbours, with competitive advantage resting on geographic proximity and low extraction costs \citep{IEA2023, Mendelevitch2019}. Newcastle benchmark prices have exhibited extreme volatility --- from \$48/tonne in 2020 to above \$400/tonne in mid-2022, normalising to \$100--140/tonne by 2024 --- illustrating the commodity price risk faced by producers and their bank lenders.

\subsection{Banking Sector and Coal Exposure}
\label{subsec:banking_exposure}

Indonesia's financial system is bank-dominated. Bank credit finances approximately 34\% of GDP, and listed banks on the IDX collectively hold over \$700 billion in total assets \citep{IMF2024Indonesia}. The sector is highly concentrated: four domestic systemically important banks (D-SIBs), namely Bank Mandiri (BMRI), Bank Rakyat Indonesia (BBRI), Bank Central Asia (BBCA), and Bank Negara Indonesia (BBNI), together account for roughly 55\% of total banking assets and 70\% of total lending. Bank Tabungan Negara (BBTN), while classified as a D-SIB for regulatory purposes, holds negligible coal exposure and is excluded from the coal-banking linkage analysis. Beyond the D-SIBs, Bank CIMB Niaga (BNGA), Bank Permata (BNLI), and Bank Syariah Indonesia (BRIS) are among the banks with material coal-sector lending.

Bank financing to the coal sector flows through project finance, corporate credit facilities (working capital, trade finance, term loans), and syndicated loans. From the exposure matrix constructed in Section~\ref{sec:data}, we identify 13 IDX-listed banks with coal exposures totalling \$6.1 billion, of which \$3.3 billion is attributed to specific bank-company pairs. BMRI holds the largest exposure (\$1.2 billion), followed by BBNI (\$975 million) and BBCA (\$441 million). The five largest bank-coal links account for over 50\% of total attributed exposure, with BMRI's lending to DSSA alone at \$586 million, implying that distress in a small number of coal companies could generate outsized losses for specific banks.

In the financial sector, Indonesian banks are heavily regulated by the Financial Services Authority (Otoritas Jasa Keuangan, OJK), which imposes a minimum capital adequacy ratio (CAR) of 8\%, aligned with Basel III standards. D-SIBs are subject to additional capital buffers, bringing the effective minimum to approximately 10.5\%. As of the latest reporting period, all 38 banks in our sample meet the minimum CAR, with D-SIBs reporting CARs ranging from 19\% to 29\%. However, these ratios are calculated on current risk weights that do not incorporate forward-looking climate transition risk, a gap that motivates our stress test.

Bank lending to the mining sector has expanded rapidly in recent years. According to OJK statistics, mining sector credit outstanding reached Rp~342 trillion (approximately \$21.6 billion) by September 2024, representing 4.5\% of total banking credit and growing at 26.7\% year-on-year. This growth is notable given the sector's historically elevated credit risk: mining sector non-performing loans peaked at 7.7\% in May 2017 during the 2015--2016 commodity price downturn, when coal prices fell below \$50/tonne. The subsequent recovery in coal prices from 2017 onward has improved asset quality and encouraged renewed bank appetite for mining sector lending, but this procyclical pattern raises the concern that banks may be extending credit precisely when transition risk is accumulating.


\subsection{Fiscal Dependence on Coal Revenue}
\label{subsec:fiscal_dependence}

The Indonesian central government derives revenue from the coal sector through three principal channels: production royalties, corporate income tax, and state-owned enterprise dividends. Together, these contributed approximately \$8.5 billion, or 4.5\% of total government revenue, in the 2022 fiscal year, though this share has fluctuated with international coal prices.


Coal royalties (PNBP) are levied at 3--7\% depending on coal type, corporate income tax applies at 22\%, and PT Bukit Asam (PTBA), the sole listed state-owned coal company, pays dividends to the government (estimated at \$130 million annually). All three revenue streams are highly procyclical: a 25\% decline in coal prices would mechanically reduce taxable income, with the effect amplified by operating leverage. Beyond the central government, coal-producing provinces in Kalimantan and South Sumatra receive revenue-sharing transfers (Dana Bagi Hasil), amplifying the political economy constraints on coal transition \citep{Heffron2018, JakobSteckel2016}.

\subsection{Regulatory Landscape and Climate Finance Policy}
\label{subsec:regulatory_landscape}

The institutional architecture governing Indonesia's climate finance response has evolved rapidly in recent years, though a significant gap remains between policy ambition and supervisory practice. Understanding this landscape is essential for contextualising the stress test results and identifying where regulatory action can most effectively reduce the financial risks we document.

Four regulatory instruments shape the transition landscape. First, \textit{OJK's Sustainable Finance Roadmap} Phase II (2021--2025) mandates ESG disclosure and calls for climate scenario analysis in bank credit risk frameworks, though implementation remains uneven. Second, \textit{Indonesia's updated NDC} (2022) commits to a 31.89\% unconditional emissions reduction by 2030, including no new coal power plant construction after 2023 and an ambition to retire 33~GW of coal capacity by 2035. Third, the \textit{Just Energy Transition Partnership} (JETP), a \$20 billion financing package launched at COP27, channels funds through an Energy Transition Mechanism for early coal plant retirement, though implementation has been slower than anticipated as of early 2026. Fourth, \textit{Bank Indonesia's Green Taxonomy} (2022) classifies economic activities by environmental impact and provides green bond eligibility frameworks.

The overall picture is one of a regulatory architecture with the right ambition but lacking the empirical underpinning to calibrate specific instruments. Capital add-ons for coal-exposed lending require quantitative estimates of default probability and loss given default under climate scenarios --- precisely the inputs our Model 2 provides. The gap between policy ambition and supervisory practice reflects a genuine information deficit that this paper helps to close.



